(1) 对 \(\eta,u\propto e^{i(kx-\omega t)}\)，线性方程化为 \((\omega-Uk)\eta_0=kh\,u_0\)、\((\omega-Uk)u_0=kg\,\eta_0\)。消元得 \((\omega-Uk)^2=ghk^2\)。记 \(c=\sqrt{gh}\)，对于正波数有 \(\omega_\pm=(U\pm c)k\)。无色散浅水近似下相速度和群速度都为 \(v_{\rm ph,\pm}=v_{\rm g,\pm}=U\pm c\)，但必须区别实验室坐标与相对于水流的传播方向。

(2) 对两支模态，\(u_0=(\omega-Uk)\eta_0/(kh)=\pm(c/h)\eta_0\)。所以顺流传播的波面升高时流体水平速度扰动同号，逆流传播时反号。与表面位移 \(a\cos(kx-\omega t)\) 对应的速度扰动振幅为 \(c a/h\)。这些关系要求 \(|\eta|\ll h,\ |u|\ll c\)，否则非线性浅水效应会使波速依赖波幅。

(3) 相对于均匀背景流，单位水平面积的二阶平均波能量为 \(\mathcal E=\langle\rho g\eta^2/2+\rho h u^2/2\rangle\)。令位移振幅为 \(a\)，势能平均 \(\rho ga^2/4\)，动能平均 \(\rho h(ca/h)^2/4=\rho ga^2/4\)，故 \(\mathcal E=\rho ga^2/2\)。若背景流速与水深处处常量，实验室系波能沿各分支的通量（相对于此背景参考能定义）为 \((U\pm c)\mathcal E\)。

(4) 逆流若 \(U<0\)，试图向 \(+x\) 传播的分支群速度为 \(U+c\)。临界 \(U=-c\) 时其实验室群速度为零；当 \(U<-c\) 时扰动被整体向负向输送。浅水模型在临界附近无色散，不能自行决定细尺度的反射和蓝移；真实有限深度色散、流速空间变化、粘性或波幅非线性可改变阻塞细节。只需判定传播方向而不能从该常系数方程宣称发生有限空间位置的自动堆积。
